\scanpage{21}
\mspage{18}

\diagram{\begin{tikzcd}[ampersand replacement=\&,row sep=small,column sep=small]
\& ? \arrow[dl,dashed,no head] \arrow[dr,dashed,no head] \& \&\\
y \arrow[dr,no head] \& \& x' \arrow[dl,no head] \arrow[dr,no head] \&\\
\& z \& \& x
\end{tikzcd}}

et il suffit de trouver $\unsure{t}$ tel que $t>\deleted{\unsure{y}}x',y$. Donc on est ramené à prouver ceci (filtration relative)

Si
\diagram{\begin{tikzcd}[ampersand replacement=\&,row sep=small,column sep=small]
\&v\\
u\arrow[ur,no head]\arrow[dr,no head]\&\\
\&w
\end{tikzcd}}
on peut compléter en
\diagram{\begin{tikzcd}[ampersand replacement=\&,row sep=small,column sep=small]
\&v\arrow[dr]\&\\
u\&\&t\\
\&w\arrow[ur]\&
\end{tikzcd}}

\deleted{\unsure{Il} manque visiblement \unsure{quelque chose}}

Mais on note qu'une relation d'ordre connexe qui est \uline{loc.} str. croissante et str. décroissante, et telle que pour \unsure{tt} élément $x$, $\exists x'_{\unsure{a}}<x<x''$, définit une structure de modèle de forme régulier connexe (on vérifie axiomes F1 à F5). Mais une telle structure n'est pas néc. filtrante croissante (\unsure{dito})

\mspage{20}

\rotatebox{90}{\reflectbox{\unsure{EXTRA STRONG}}}
\rotatebox{90}{\reflectbox{\unsure{EXTRA STRONG}}}
